Este trabajo de investigación surge con el propósito fundamental de aportar evidencia empírica y valor analítico a la industria del desarrollo de videojuegos. En las etapas iniciales de concepción del proyecto de graduación existieron diversas propuestas, todas unidas por el interés común y la vocación de diseñar un videojuego. Sin embargo, el surgimiento y la adopción de herramientas de inteligencia artificial en la industria marcaron el rumbo definitivo del estudio. El debate contemporáneo generado por títulos comerciales que, a pesar de su alta calidad técnica y jugabilidad, experimentaron una pérdida sustancial de credibilidad debido a la divulgación sobre el uso de IA en su proceso creativo, motivó la necesidad de investigar formalmente este fenómeno de percepción.

Deseo expresar un sincero agradecimiento a Nicolle Gordillo, coautora de este trabajo experimental colaborativo, por la elaboración de la totalidad del arte en \textit{pixel art} utilizado en los prototipos. De igual forma, reconozco el valioso aporte de Héctor Penedo, a pesar de no formar parte del equipo de investigación formal, contribuyó con la composición y producción de la música original del videojuego.

Este trabajo no habría alcanzado su madurez técnica y metodológica sin el acompañamiento del profesor Dennis Aldana, catedrático de la Universidad del Valle de Guatemala, cuya asesoría inicial fue clave para estructurar la idea de investigación. Asimismo, agradezco profundamente al Ing. Alexander Bolaños por su orientación constante, rigor y guía académica a lo largo del proceso de elaboración de este trabajo de graduación.

Por último, y de manera sumamente especial, extiendo mi gratitud a los 60 participantes que dedicaron su tiempo a evaluar los prototipos e interactuar con el entorno experimental. Su disposición y valiosa retroalimentación constituyeron el pilar fundamental que hizo posible la ejecución y el éxito de esta investigación. 